% =====================================================================================
\section{Discussion}
\label{sec:discussion}
% =====================================================================================

\subsection{Comparison with published devices}
\label{sec:compare}

Our results are simulated and the published ones are measured unless marked otherwise. Steering power is compared at
a fixed angle (\ang{10}) and scaled to $N=33$, because a wrapped array draws nearly the same power at any angle and
power per degree would favour it at large angles.

\paragraph{Few-wire steering.} The directly comparable designs are those that also steer with a few signals
(\cref{tab:fewwire}). Measured cascaded arrays at our pitch need about \qtyrange{12}{17}{\mWdeg} when scaled to
\num{33} emitters; the \stack{slot} and \stack{vacuum} designs need \num{1.38} and \qty{0.53}{\mWdeg},
\numrange{10}{30} times less. The cascaded arrays are much faster (\qty{100}{\kilo\hertz} bandwidth~\cite{yaacobi2014}
against our \qtyrange{50}{190}{\hertz} with plain steps), but their heated stages are micrometres long and run near
\qty{500}{\celsius}, and their series feed accumulates loss and error along the bus.

\begin{table}[htbp]
  \centering
  \caption{Steering with few control signals. Power per degree of the published arrays is scaled to $N=33$ from the
  reported totals (our estimate).}
  \label{tab:fewwire}
  \footnotesize
  \setlength{\tabcolsep}{4pt}
  \begin{tabularx}{\linewidth}{@{}P{0.18\linewidth}LcP{0.1\linewidth}P{0.1\linewidth}cP{0.11\linewidth}@{}}
    \toprule
    Work & Ramp made by & Wires & Array & Range & \unit{\mWdeg} & Basis \\
    \midrule
    This work, \stack{vacuum} & heat flow across a slotted membrane & 1 & $33\times\qty{2}{\micro\metre}$ &
        \ang{22.8} one way & 0.53 & simulated \\
    This work, \stack{slot} & same, in air & 3 & same & \ang{22.8} one way & 1.38 & simulated \\
    This work, \stack{undercut} & membrane anchored at both walls & 4 & same & \ang{\pm 22.8} & 5.24 & simulated \\
    \addlinespace
    Yaacobi~\cite{yaacobi2014} & cascaded heaters on one bus & 1 & $16\times\qty{2}{\micro\metre}$ & \ang{51} &
        $\approx12$ & measured, scaled \\
    Poulton~\cite{poulton2016} & cascade in three groups & 3 & $50\times\qty{2}{\micro\metre}$ & \ang{46} &
        $\approx17$ & measured, scaled \\
    Van Acoleyen~\cite{vanacoleyen2009} & one triangular electrode & 1 & $16\times\qty{2}{\micro\metre}$ & \ang{2.3}
        & $\propto N^2$ & measured \\
    Hosseini~\cite{hosseini2025} & meander overlap (SiN, \qty{525}{\nano\metre}) & 4 & 16 & \ang{\pm 5.1} & --- &
        measured \\
    \bottomrule
  \end{tabularx}
\end{table}

\paragraph{One heater per emitter.} Against arrays of the best individual phase shifters the result depends on the
device (\cref{sec:power}): at \ang{10} the \stack{slot} design equals \num{33} single-pass suspended
heaters~\cite{fang2011} and the \stack{vacuum} design needs \num{3.6} times more than \num{33} folded record
shifters~\cite{lu2015}. These arrays have not been built; the record devices are large (a \qty{6}{\milli\metre}
folded arm~\cite{lu2015}) and would need a fan-in and \num{33} drivers. Measured thermo-optic OPAs draw
\qtyrange{3.7}{9.4}{\milli\watt} per element~\cite{miller2020,phare2018}, against our
\qtyrange{0.16}{0.42}{\milli\watt}, but their totals include the power used to cancel fabrication errors, which our
design replaces with a trim.

\paragraph{Lenses.} The two measured few-heater lenses found cost about \qty{40}{\milli\watt\per\rad} (shaped heaters
on a silicon slab metalens~\cite{liu2023}) and \qtyrange{70}{110}{\milli\watt\per\rad} (meanders in silicon
nitride~\cite{hosseini2025}), converted by us from the reported focal shifts and good to about a factor of two. Ours,
at about \qty{1}{\milli\watt\per\rad}, is \numrange{30}{100} times lower. Most of this difference is due to the
platform: those devices sit on solid substrates, and silicon nitride's thermo-optic coefficient is about seven times
smaller than silicon's. The metalens is about ten times faster.

\Cref{tab:scoreboard} summarises the comparison.

\begin{table}[htbp]
  \centering
  \caption{Summary of the comparison. Ours: two-dimensional simulation. Theirs: measured, or, where marked $\dagger$,
  our estimate for \num{33} copies of a measured device.}
  \label{tab:scoreboard}
  \footnotesize
  \begin{tabularx}{\linewidth}{@{}P{0.17\linewidth}P{0.25\linewidth}LP{0.2\linewidth}@{}}
    \toprule
    Metric & This work & Best published & Outcome \\
    \midrule
    Few-wire steering power & \qtyrange{0.53}{1.38}{\mWdeg}, 1--3 wires & \qtyrange{12}{17}{\mWdeg}, 1--3
        wires~\cite{yaacobi2014,poulton2016} & $10$--$30\times$ lower \\
    Power at \ang{10} & \qty{5.3}{\milli\watt} (\stack{vacuum}), \qty{13.8}{\milli\watt} (\stack{slot}) &
        \qty{1.5}{\milli\watt}$^\dagger$~\cite{lu2015}; \qty{14.4}{\milli\watt}$^\dagger$~\cite{fang2011} &
        $3.6\times$ higher than the record; equal to single-pass \\
    Tunable lens & \qtyrange{0.9}{1.2}{\milli\watt\per\rad}, 1--5 wires & \qtyrange{40}{110}{\milli\watt\per\rad}
        (est.)~\cite{liu2023,hosseini2025} & $30$--$100\times$ lower \\
    Speed, \ang{0}$\to$\ang{10} & \qtyrange{2}{9}{\milli\second} plain; \qtyrange{0.16}{1.4}{\milli\second}
        with pre-emphasis & \qty{2}{\micro\second}~\cite{jacques2019} to \qty{144}{\micro\second}~\cite{fang2011}
        per device & slower; power $\times$ time $1.2$--$16\times$ higher with pre-emphasis \\
    Field of view & \ang{\pm 22.8} (one way for \stack{slot}, \stack{vacuum}) & \ang{51} at the same
        pitch~\cite{yaacobi2014} & comparable within pitch class \\
    Calibration & one trim, then a linear fit and a gain curve & a search per beam~\cite{lian2024} & simpler, given
        the trim \\
    Independent errors & not correctable; trim required & fully corrected & worse \\
    Patterns & ramp and dome only & arbitrary & worse \\
    \bottomrule
  \end{tabularx}
\end{table}

\subsection{Agreement with real devices}
\label{sec:fidelity}

The single-heater versions of our thermal model agree with measured devices. Our \stack{standard} heater has
$\Ppi=\qty{19.7}{\milli\watt}$, against \qtyrange{22.8}{29.2}{\milli\watt} measured for TiN and doped-silicon
heaters~\cite{jacques2019}; our \stack{undercut} heater has \qty{1.35}{\milli\watt} and \qty{1.6}{\milli\second},
against \qty{1.07}{\milli\watt} and \qtyrange{1.3}{1.9}{\milli\second} for a foundry undercut~\cite{parsons2025}.
Published comparisons of simulated and measured thermo-optic devices agree to about $\pm\qty{15}{\percent}$ for
conventional heaters and within a factor of \numrange{0.67}{1.5} for undercut structures with idealised geometry. A
real back-end stack would also add oxide above the heater, increasing $\kt$ and the power by \numrange{1.5}{1.9}
times. A fabricated \stack{slot} design would therefore need roughly \qtyrange{9}{40}{\milli\watt} at \ang{10}, and a
\stack{vacuum} design \qtyrange{3.5}{15}{\milli\watt} (our estimate), still several times below the measured few-wire
arrays. The structure itself, a single \qty{80}{\micro\metre} membrane under a whole array, has not been fabricated;
published undercuts isolate single devices with undercuts of order \qty{10}{\micro\metre}~\cite{parsons2025}.
The \stack{vacuum} design also needs a sealed package. Gas conduction across the \qty{15}{\micro\metre} cavity
falls only once the mean free path approaches the gap, and a simple kinetic estimate puts the required pressure
near \qty{100}{\pascal}; no thermo-optic device has, to our knowledge, been measured in vacuum.

\subsection{Limitations}
\label{sec:limitations}

\Cref{tab:limits} lists the limitations, the claims they affect and how each could be resolved. The most important is
that all thermal results are two-dimensional: a finite heater loses some heat through the cavity ends (an estimated
\qtyrange{5}{8}{\percent} more power for \qty{1}{\milli\metre} heaters), and heat conducted along the waveguides could
reach the gratings. In addition, the \stack{slot} and \stack{vacuum} designs steer in one direction only, and the
design does not scale unchanged to large apertures: at \num{128} emitters the hot end would reach about
\qty{250}{\kelvin} at \ang{10}, and settling would be about \num{16} times slower.

\begin{table}[htbp]
  \centering
  \caption{Limitations.}
  \label{tab:limits}
  \footnotesize
  \begin{tabularx}{\linewidth}{@{}P{0.24\linewidth}P{0.22\linewidth}LP{0.18\linewidth}@{}}
    \toprule
    Limitation & Affects & Expected size & Resolved by \\
    \midrule
    2D model; no heater or cavity ends & all powers; grating temperature & $+$\qtyrange{5}{8}{\percent} power (est.);
        grating heating unknown & 3D thermal model \\
    No fabricated device & absolute values & $\times0.67$--$1.5$; $\times1.5$--$1.9$ for a real stack & test chip \\
    Idealised release and \qty{2}{\micro\metre} slot & \stack{slot} power & a wider slot and deeper cavity help &
        realistic-stack model \\
    Mechanics not modelled & membrane feasibility & buckling margin, slot-edge curl unknown & thermo-mechanical
        model \\
    Vacuum operation unprecedented & \stack{vacuum} results & --- & measurement against pressure \\
    Emitters identical and independent & far-field model & small ($L_c=\qty{2.2}{\metre}$) & 3D full-wave model \\
    Static trim assumed & beam quality beyond small angles & $S\approx0.78$ untrimmed at \ang{10} & trim shown on
        an OPA \\
    Silicon's $\mathrm{d}n/\mathrm{d}T$ growth unconfirmed & sign of pointing error & $\pm\ang{0.25}$ at \ang{10}
        before gain correction & coefficients of Ref.~\cite{cocorullo1999} \\
    Kernel-shape errors not modelled & calibration & expected to rescale the ramp & sensitivity runs \\
    One-way steering (\stack{slot}, \stack{vacuum}) & field of view & half the field, or twice the power &
        two-sided layout \\
    \bottomrule
  \end{tabularx}
\end{table}

\subsection{Applications and further work}

The design is a low-wire, low-power, slow actuator for smooth phase profiles over a small aperture. It suits tasks
that need a ramp or a parabola of variable strength at rates below about \qty{100}{\hertz}, such as tip-tilt
correction behind a telescope in ground-level free-space links, where loop bandwidths of \qtyrange{1}{70}{\hertz} are
required~\cite{frost2024}, coarse beam acquisition and variable focus. It does not suit LiDAR scanning, which needs
\qtyrange{10}{100}{\kilo\hertz}, or adaptive optics, which needs arbitrary phase profiles.

The next steps, all specified as COMSOL models, are: a three-dimensional model of a finite heater and cavity including
the gratings; a model with a realistic back-end stack, a wider release slot and an isotropic cavity; the
thermo-mechanical bow of the slotted membrane; kernel sensitivities to fabrication variations; and a two-sided slot
layout.
